\documentclass[11pt,a4paper]{article}
\usepackage{fontspec}
\setmainfont{Noto Serif}
\setsansfont{Noto Sans}
\usepackage{babel}
\babelprovide[main,import]{russian}
\usepackage{amsmath,amssymb,mathtools,booktabs,geometry,fancyhdr,microtype,xcolor,hyperref}
\geometry{left=21mm,right=21mm,top=21mm,bottom=21mm}
\definecolor{ink}{HTML}{173D57}
\hypersetup{colorlinks=true,urlcolor=ink,linkcolor=ink,pdftitle={BGS-F. Убывание знаковой резольвенты}}
\pagestyle{fancy}\fancyhf{}\lhead{\small BGS-F}\rhead{\small Убывание знаковой резольвенты}\cfoot{\thepage}
\setlength{\headheight}{15pt}\setlength{\parindent}{1em}\setlength{\parskip}{2.5pt}
\setlength{\emergencystretch}{3em}\tolerance=3000
\newcommand{\Tr}{\operatorname{Tr}}
\newcommand{\EE}{\mathbb E}
\newcommand{\R}{\mathbb R}
\newcommand{\Z}{\mathbb Z}
\newcommand{\C}{\mathbb C}
\newcommand{\si}{\mathrm{Sine}_2}
\newcommand{\eps}{\varepsilon}
\newcommand{\sinc}{\operatorname{sinc}}
\newcommand{\sect}[1]{\section*{\sffamily\color{ink}#1}}
\newcommand{\subt}[1]{\par\medskip\noindent\textbf{#1.}\quad}
\newcommand{\T}{\mathbb T}
\newcommand{\dist}{\operatorname{dist}}
\newcommand{\Ad}{\operatorname{Ad}}
\newcommand{\op}{\mathrm{op}}
\newcommand{\HS}{\mathrm{HS}}
\newcommand{\Zp}{\mathbb Z_7}
\begin{document}
\noindent{\small\sffamily\color{ink} МАТЕМАТИЧЕСКАЯ ФИЗИКА\hfill 20 НОЯБРЯ 2026}
\begin{center}{\LARGE\sffamily Убывание знаковой резольвенты}\\[4pt]
{\large Локальная теорема с потерей \(\log\log N\)}
\end{center}
Пусть \(N=7^k\), \(m=3N/7\), \(U_N(t)=e^{itA_N}M_N\). Оператор \(A_N\) и линейный пропагатор \(M_N\) взяты из исходной модели BGS-F. Физическая сила равна \(u_N=t/N\).

\subt{Новый предел}
Если
\[
                    u_N=o\!\left(\frac{N^{-3/2}}{\log\log N}\right),
\]
локальный спектр сходится к случайно сдвинутой решётке шага \(7/3\) с кратностями \(2(1+\nu_7(B+j))\), где \(B\) равномерен в \(\Zp\). Для каждого \(T<\infty\) обычные следы удовлетворяют
\[
    \sup_{|r|\le TN}\frac{|\Tr U_N(t)^r-\Tr M_N^r|}{N}\longrightarrow0 .
\]
В частности, при \(r=m\) модуль нормированного следа стремится к единице. Для некоторой ненулевой константы \(c\) закон \(\si\) также исключается при \(|u_N|\le cN^{-3/2}/\log\log N\).

\subt{Количественная оценка}
Фиксируем \(R\in(0,1]\) и проводим окружности радиуса \(R/m\) вокруг линейных собственных значений. При достаточно малом \(|t|\sqrt N\) линейная кратность сохраняется внутри каждой окружности, кроме доли не более
\[
    C_R\left[t^4N^2+\min\left(t^2N(1+\log k)^2,\ t^4N^2k\right)\right].
\]
Основное улучшение состоит в оценке среднего квадрата нормы квадратичного знакового блока величиной \(C_Rt^4N^2\), тогда как предыдущий счёт давал \(C_Rt^4N^2k\). После разбиения замкнутого слова из четырёх меток на арифметические классы коэффициенты знаковой резольвенты суммируются по каждому классу с убыванием.

Для полного парного теста остаётся справедливым ранее полученный предел \((0.08443,0.08656)\) при \(u_N=o(N^{-8/5}(\log N)^{-2/5})\). Новая оценка расширяет область локальной сходимости и контроля обычных следов. Она использует убывающую силу; исходная гипотеза при \(u_N=10^{-80}\) остаётся открытой.
\newpage
\sect{1. Убывание на произвольной прогрессии}
Как в предшествующих частях, \(\lambda_j=\lambda_0e^{2\pi i j/m}\), \(E_j\) --- ортогональные линейные спектральные проекторы. Представитель минимального модуля \(d\bmod m\) имеет \(|d|\le D=(m-1)/2\). Пусть \(L=\lfloor k^{16}\rfloor\). Все оценки далее относятся к достаточно большим \(k\), когда \(1<L<D\). Обозначим
\[
 w_0=R^{-1},\qquad w_d=(2m\sin(\pi |d|/m)-R)^{-1}\ (d\ne0),\qquad
 w^{\rm n}_d=w_d\mathbf1_{|d|\le L},\quad
 w^{\rm f}_d=w_d\mathbf1_{|d|>L}.
\]
Для неотрицательного веса \(g\) положим \(Z_j(g)=\sum_i\sqrt{g_{i-j}}E_i\), \(Z_{\rm n}=Z_j(w^{\rm n})\), \(Z_{\rm f}=Z_j(w^{\rm f})\), \(Z=Z_j(w)\). Дальняя резольвента с фазой в центре окружности есть \(mF_j\), где
\[
             F_j=\sum_i f_{i-j}E_i,\qquad
       f_d=\frac{\mathbf1_{|d|>L}}{m(1-e^{-2\pi i d/m})}.
\]
Преобразование \(\widehat f(h)=\sum_df_d e^{-2\pi i hd/m}\) взято без нормировки.

Для \(1\le d\le D\) имеем
\[
  \frac1{m(1-e^{-2\pi i d/m})}=\frac1{2m}-\frac{i}{2m}\cot(\pi d/m).
\]
Вместе с отрицательным \(d\) нечётная часть даёт усечённый синусный гармонический ряд, с ошибкой \(O(1)\) по сумме модулей, поскольку \(\frac1{2m}\cot(\pi d/m)=\frac1{2\pi d}+O(d/m^2)\). Синусная сумма \(\sum_{a\le d\le b}\sin(\theta d)/d\) ограничена равномерно: до \(d=1/|\theta|\) используем \(|\sin(\theta d)|\le|\theta|d\), дальше --- суммирование по частям геометрической суммы. Поэтому \(|\widehat f(h)|\le C\) для всех \(h\). Дополнительно на каждой из двух полуосей \(f_d\) при \(|d|>L\) имеет вариацию и максимальный модуль \(O(1/L)\). Суммирование по частям даёт
\begin{equation}
       |\widehat f(0)|\le C,\qquad
       |\widehat f(h)|\le C\min\{1,m/(L|h|_m)\}\quad(h\ne0).
                                                               \tag{1}
\end{equation}
То же верно для преобразования \(\bar f\).

Если \(q\mid m\), то точки произвольного класса \(h\equiv a\bmod q\) равноотстоят. По мере удаления от нуля их расстояния, кроме ближайших, возрастают хотя бы как \(q l/2\), \(l\ge1\). Из (1) и \(\sum_{l\ge1}\min(1,b^2/l^2)\le Cb\) получаем
\begin{equation}
                \sum_{h\equiv a\bmod q}|\widehat f(h)|^2
                          \le C(1+m/(qL)).
                                                               \tag{2}
\end{equation}
Граница равномерна по сдвигу \(a\), включая класс, проходящий через частоту \(0\). Именно эта равномерность нужна при завершении вейлева слова.
\newpage
\sect{2. Замкнутое слово и прогрессии времён}
Достаточно оценить \(Z_{\rm n}T(v)F_jT(w)Z_{\rm n}\) для каждой пары исходных вейлевых меток. Квадрат его нормы Гильберта--Шмидта раскладывается по словам с метками \(v,w,-w,-v\) и весами \(F_j,Z_{\rm n}^2,F_j^*,Z_{\rm n}^2\) между ними. Как в части о четвёртом моменте, после усреднения по \(j\) четыре частоты имеют сумму \(0\bmod m\). Обозначим накопленные времена меток через \(0,s_1,s_2,h\). Продукт весов по модулю имеет вид
\[
 |\widehat f(s_1)\widehat{\bar f}(h-s_2)
                   \widehat w^{\rm n}(s_2-s_1)\widehat w^{\rm n}(-h)|.
\]
Используем \(2|ab|\le |a|^2+|b|^2\) отдельно к двум ближним весам и к двум дальним. Если остаётся \(|\widehat w^{\rm n}(s_2-s_1)|^2\), циклически начинаем слово с \(-w\) и снимаем общее время \(s_2\). Тогда этот вес зависит от последнего времени, а метки \(v,w\) заменяются на \(-w,-v\). Таким образом, в каждом слагаемом можно фиксировать \(h\) с внешним весом \(|\widehat w^{\rm n}(h)|^2\), а во внутренней сумме оставить один из весов \(|\widehat f(s_1)|^2\), \(|\widehat{\bar f}(h-s_2)|^2\).

В расщеплённых координатах действие \(S\) равно \((\gamma,\gamma^{-1})\) с \(\gamma=3+2\rho,\ \rho^2=2,\ \rho\equiv3\bmod7\). После снятия общей степени семи с \(v,w\) модуль есть \(7^K\), \(K=k\) или \(k-1\). Запишем \(v=7^{\delta_0}(a,b)\), \(w=7^{\delta_1}(c,d)\), где координаты \(a,b,c,d\) обратимы, \(\delta_i\in\{0,1\}\) и \(\min(\delta_0,\delta_1)=0\). При \(X=\gamma^h,Y=\gamma^{s_1},Z_0=\gamma^{s_2}\) замыкание эквивалентно
\begin{equation}
  7^{\delta_1}c(Y-Z_0)=7^{\delta_0}a(X-1),\qquad
  7^{\delta_0}(X-1)(YZ_0-\eta X)=0,\quad
                  \eta=\frac{ad}{bc}.
                                                               \tag{3}
\end{equation}
Эта формула получается домножением уравнения для обратных координат на \(XYZ_0\) и подстановкой первого уравнения. След каждого замкнутого слова по модулю равен \(N\).

Пусть \(s=\min(K,\nu_7(X-1))\). При \(s+\delta_0<\delta_1\) завершений нет. При \(\delta_1\le s+\delta_0<K\), положив
\(C_0=(a/c)7^{\delta_0-\delta_1}(X-1)\), \(K'=K-s-\delta_0\), получаем
\[
 Z_0=Y-C_0\bmod7^{K-\delta_1},\quad
 Y^2-C_0Y-\eta X=0\bmod7^{K'},\quad
 Z_0^2+C_0 Z_0-\eta X=0\bmod7^{K'}.
\]
Время, приводящее к каждому заданному корню по модулю \(7^{K'}\), лежит в одном классе по модулю \(q=3\cdot 7^{K'-1}\) или отсутствует. Это следует из точного порядка \(\gamma\bmod 7^{K'}\). Если фиксировать такое время \(s_1\), то \(s_2\) имеет не более \(7\) вариантов по первому равенству, и наоборот: число есть не более \(m/(3\cdot7^{K-\delta_1-1})\le7\). Таким образом, для каждого корня по (2) взвешенная сумма завершений не больше
\begin{equation}
                        C(1+7^s/L).
                                                               \tag{4}
\end{equation}
Для второго внутреннего веса отражение и сдвиг \(h-s_2\) просто меняют класс прогрессии.
\newpage
\sect{3. Суммирование по порядку делимости}
\subt{Невырожденное завершение}
При \(s+\delta_0>\delta_1\) и \(s+\delta_0<K\) коэффициент \(C_0\) делится на \(7\). Дискриминант \(C_0^2+4\eta X\) обратим, поэтому допустимы не более двух корней. Время \(h\ne0\), а при \(s\ge1\) оно кратно \(3\cdot7^{s-1}\). Для ближнего веса
\[
       |\widehat w^{\rm n}(h)|\le C_R[1+\log(m/(1+|h|_m))],
                       \quad h\ne0.
\]
Гармоническое сравнение на каждой подгруппе даёт \(\sum_{h\ne0,\ \nu_7(X-1)\ge s}|\widehat w^{\rm n}(h)|^2\le C_Rm7^{-s}\). При \(s=0\) та же оценка верна по всем ненулевым \(h\). Умножая на (4) и суммируя по \(s\), получаем
\[
                  C_R m\sum_{s=0}^{K-1}(7^{-s}+L^{-1})
                                     \le C_R m(1+K/L).
\]
Различие от предыдущей оценки принципиально: множитель \(7^s\) у числа завершений заменяется на \(1+7^s/L\) до суммирования уровней делимости.

\subt{Исключительный дискриминант}
Пусть \(s+\delta_0=\delta_1\). Тогда возможны только \((s,\delta_0,\delta_1)=(0,0,0)\) и \((1,0,1)\). В первом случае дискриминант по \(X\) есть \((a/c)^2(X-1)^2+4\eta X\). Его корни простые по модулю 7, кроме \(\eta=(a/c)^2\bmod7\), когда потенциальный корень \(-1\) не входит в группу \(\langle\gamma\rangle\bmod7\). Во втором случае, при \(X=1+7X_0\), дискриминант по \(X_0\) также имеет простые корни по модулю 7, так как \(\eta\) обратима. Поднятие корней даёт \(O(m7^{-b_0})\) времён, где дискриминант делится на \(7^{b_0}\), при \(b_0\le K'\). Во втором случае для такого подсчёта нужен \(X\bmod7^{b_0+1}\). Число завершений при порядке дискриминанта \(b_0\), усечённом на \(K'\), не больше \(C7^{b_0/2}\), включая поднятия: \(s\) здесь равен 0 или 1. Пусть \(B(h)\) обозначает это верхнее ограничение, нулевое для других \(h\). Тогда \(\sum_h B(h)^{3/2}=O(m)\) благодаря \(\sum_{b_0\ge0}7^{-b_0/4}<\infty\). По Гёльдеру
\[
 \sum_{h\ne0} |\widehat w^{\rm n}(h)|^2 B(h)
 \le
 \left(\sum_{h\ne0}|\widehat w^{\rm n}(h)|^6\right)^{1/3}
 \left(\sum_h B(h)^{3/2}\right)^{2/3}
 \le C_R m .
\]
Первая сумма \(O_R(m)\) по интегрируемости \((1+\log(1/x))^6\) на \((0,1)\). Дальние веса здесь равномерно ограничены по (1).

\subt{Внешние метки погашаются}
При \(s+\delta_0\ge K\) имеется \(O(1)\) внешних времён \(h\), поскольку модуль совместимости не меньше \(7^{K-1}\). Первое равенство (3) даёт \(Y=Z_0\bmod7^{K-\delta_1}\). Для каждого времени одной из внутренних меток допустимо лишь \(O(1)\) времён другой. По (2) с \(q=1\) внутренняя сумма \(\le C(1+m/L)\). Внешний вес не больше \(C_R(1+\log(2L))^2\), что даёт \(O_R(m)\) при \(L=\lfloor k^{16}\rfloor\).

Итак, взвешенное число замкнутых слов \(\le C_R m\). Четыре знаменателя \(m^{-1}\) из разложения весов сокращаются с \(m^4\) в квадрате нормы. Используя \(N\) для следа и треугольник по конечным исходным коэффициентам в \(L^2(j;\mathrm{HS})\), получаем
\begin{equation}
       \EE_j\|t^2m^2 Z_{\rm n} A_N F_j A_N Z_{\rm n}\|_{\HS}^2
                             \le C_R t^4 N^2 .
                                                               \tag{5}
\end{equation}
\newpage
\sect{4. Количественный подсчёт спектральных кластеров}
Для \(Z_{\rm n},Z_{\rm f},Z\) используем предыдущие четвёртые оценки
\[
\begin{aligned}
 \EE_j\|tm Z_{\rm n}A_N Z_{\rm n}\|_{S_4}^4&\le C_Rt^4N^2k,&
 \EE_j\|tm Z_{\rm f}A_N Z_{\rm f}\|_{S_4}^4&\le C_Rt^4N^2 k^3/L,\\
 \EE_j\|tm Z A_N Z\|_{S_4}^4&\le C_R t^4N^2 k^2 .
\end{aligned}
\]
Дополнительно в ближней части достаточно второй нормы. Из равномерной энергии между узлами \(W_h=\sum_i\|E_i A_NE_{i+h}\|_{\HS}^2\le C N/m\) и \(\sum_d w^{\rm n}_d=O_R(1+\log(2L))\) следует
\[
        \EE_j\|tm Z_{\rm n} A_N Z_{\rm n}\|_{\HS}^2
                                    \le C_R t^2N(1+\log(2L))^2 .
\]
Усреднение по \(j\) даёт \(t^2m\sum_h W_h\sum_d w^{\rm n}_d w^{\rm n}_{d+h}\), откуда и получается граница. Оценку \(W_h\) даёт разложение по орбитам вейлевых меток, подробно в части о резольвенте.

Выбираем хорошие узлы условиями
\[
 \begin{gathered}
 \|tmZ_{\rm n} A_N Z_{\rm n}\|_{\op}\le1/4,\qquad
 \|t^2m^2Z_{\rm n}A_NF_jA_NZ_{\rm n}\|_{\op}\le1/4,\\
 \|tm Z_{\rm f} A_N Z_{\rm f}\|_{\op}\le k^{-2},\qquad
 \|tm Z_{\rm n} A_N Z_{\rm f}\|_{\op}\le k^{3/4}.
 \end{gathered}
\]
Для первого условия применяем лучшее из двух средних, для второго (5), для остальных дальнюю и полную четвёртые оценки. По Маркову доля остальных узлов ограничена
\begin{align}
   \beta\ &\le C_R\!\left[\min\{t^2N(1+\log(2L))^2,t^4N^2k\}
                     +t^4N^2(1+k^{11}/L+k^{-1})\right] \notag\\
   &\le C_R'\!\left[t^4N^2+\min\{t^2N(1+\log k)^2,t^4N^2k\}\right].
                                                               \tag{6}
\end{align}

Проверим, что на хорошем узле действительно сохраняется кратность. На \(\Gamma_j\) берём \(D_z=(M_N-zI)^{-1/2}\), \(B(z)=D_z(e^{itA_N}-I)M_ND_z\), \(B^{(1)}(z)=itD_z A_N M_N D_z\). Линейные блоки покомпонентно ограничены соответствующими весовыми нормами, поскольку \(|D_z|\le\sqrt m Z\) в линейном спектральном базисе. По \(\|e^{itA_N}-I-itA_N\|_{\op}\le \tfrac12 t^2\|A_N\|_{\op}^2\) остатки в ближнем, дальнем и смешанных блоках ограничены \(C_Rt^2 N\), \(C_Rt^2N/L\), \(C_Rt^2N/\sqrt L\). При \(|t|\sqrt N\le c_R\) с достаточно малым фиксированным \(c_R\) ближний остаток не превышает \(1/8\).

В \(I+sB(z)\), \(s\in[0,1]\), устраняем дальний блок через дополнение Шура. Его обратная матрица равна \(I+O(k^{-2})\), так что замена её на \(I\) в смешанном произведении даёт \(O_R(k^{-2}k^{3/2})=o(1)\). Замена самих смешанных блоков на линейные стоит \(O_R(k^{3/4}/\sqrt L)=o(1)\). В оставшемся квадратичном блоке \(B_{\rm nf}^{(1)}B_{\rm fn}^{(1)}\) замена \(P_{\rm f}M_N(M_N-zI)^{-1}\) на центральный \(mF_j\) стоит \(O_R(k^{3/2}/L)=o(1)\) по тождеству резольвент: разность покомпонентно \(\le m w^{\rm f}_{i-j}R/(4L)\). Полученный блок контролируется вторым условием хорошего узла. В итоге дополнение Шура отличается от \(I\) по норме не больше \(1/4+1/8+1/4+o(1)<1\), равномерно на контуре и по \(s\). Принцип аргумента для \(M_N+s(e^{itA_N}-I)M_N\) даёт сохранение числа уровней внутри \(\Gamma_j\).
\newpage
\sect{5. Переход к локальному спектру}
Кратности линейных фаз даны формулой
\(d_j=2\sum_{a=0}^{k-1}\mathbf1_{7^a\mid j}+\mathbf1_{j=0}\).
В каждой хорошей окружности сопоставим все \(d_j\) уровней с её центром. Для доли оставшихся фаз \(\epsilon\) получаем
\[
  \epsilon\le(6/7)\sum_{a=0}^{k-1}\min(\beta,7^{-a})+1/N,
         \qquad \epsilon\le Ck\beta .
\]
В частности, если \(\beta\to0\), то \(\epsilon\to0\) при любой фиксированной ширине \(R>0\). Сопоставленные уровни отличаются после развёртки на \(O(R)\).

При \(t=o(N^{-1/2}/\log k)\) по (6) \(\beta\to0\). При общем равномерном начале развёртки и \(f\in C_c^1(\R)\) вклад сопоставленных уровней в разность линейных статистик в среднем не больше \(C R\|f'\|_1\) по интегрированию периодической копии \(f\), оставшихся --- не больше \(2\epsilon\|f\|_1\). Сначала \(k\to\infty\), затем \(R\downarrow0\) переносят весь локальный закон; теснота следует из единичной интенсивности. Формула \(d_j\) описывает предел линейного закона: случайно сдвинутая решётка шага \(N/m=7/3\) с общим \(\Zp\)-индексом \(B\), кратности \(2(1+\nu_7(B+j))\); добавочный уровень при \(j=0\) исчезает по вероятности. Этот процесс отличается от \(\si\), для которого нет кратных точек.

На временах \(|r|\le TN\) сопоставление даёт
\[
        N^{-1}|\Tr U_N(t)^r-\Tr M_N^r|\le C_T R+2\epsilon,
\]
значит, предел совпадает равномерно. Если лишь \(|t|\le c N^{-1/2}/\log k\), при фиксированном \(R\) доля пропущенных узлов ограничена \(C_Rc^2+o(1)\). Выбираем малый \(R\), затем малый положительный \(c\). На положительной доле узлов сохраняются две или больше близкие фазы, и ограниченный тест на близкие пары из части о резольвенте отделяет закон от \(\si\). Верхняя граница для вероятности близкой пары у \(\si\) кубически убывает с шириной; для сохранённых узлов вероятность попасть в фиксированное локальное окно имеет положительную нижнюю границу.

Для парной статистики используется вторая граница для \(\epsilon\), и стандартное сопоставление из предшествующих частей даёт для теста Шварца \(\psi\)
\[
   |P_\psi(U_N(t))-P_\psi(M_N)|\le C_\psi[R+\epsilon(1+N|t|+k)] .
\]
Вычитание по сопоставленным парам контролируется \(\sum_j d_j^2=O(N)\) и интегрируемой мажорантой производной \(\psi\). Для прочих пар число уровней в единичном интервале не больше \(C(1+N|t|+k)\) по скорости фаз и линейным кратностям. Из четвёртого варианта в (6) сохраняется предел в области \(t=o(N^{-3/5}k^{-2/5})\). Это относится и к исходному \(\Delta_{2,N}(b_*)\). Переход к \(u_N=t/N\) даёт формулировки первой страницы.

\subt{Проверка и граница}
Целочисленный код проверяет обе квадратные формулы и прогрессии времён при \(N=49,343,2401\). Отдельный численный тест проверяет убывание коэффициентов и суммы по сдвинутым прогрессиям. Оценка (6) использует ограничение на \(|t|\sqrt N\), тогда как при фиксированном \(10^{-80}\) оно растёт как \(N^{3/2}/10^{80}\).
\end{document}
